%; whizzy paragraph -pdf xpdf -latex ./whizzypdfptex.sh
%; whizzy-paragraph "^\\\\begin{frame}"
% latex beamer presentation.
% platex, latex-beamer でコンパイルすることを想定。 

%     Tokyo Debian Meeting resources
%     Copyright (C) 2009 Junichi Uekawa
%     Copyright (C) 2009 Nobuhiro Iwamatsu

%     This program is free software; you can redistribute it and/or modify
%     it under the terms of the GNU General Public License as published by
%     the Free Software Foundation; either version 2 of the License, or
%     (at your option) any later version.

%     This program is distributed in the hope that it will be useful,
%     but WITHOUT ANY WARRANTY; without even the implied warreanty of
%     MERCHANTABILITY or FITNESS FOR A PARTICULAR PURPOSE.  See the
%     GNU General Public License for more details.

%     You should have received a copy of the GNU General Public License
%     along with this program; if not, write to the Free Software
%     Foundation, Inc., 51 Franklin St, Fifth Floor, Boston, MA  02110-1301 USA

\documentclass[cjk,dvipdfmx,12pt]{beamer}
\usetheme{Tokyo}
\usepackage{monthlypresentation}

%  preview (shell-command (concat "evince " (replace-regexp-in-string "tex$" "pdf"(buffer-file-name)) "&")) 
%  presentation (shell-command (concat "xpdf -fullscreen " (replace-regexp-in-string "tex$" "pdf"(buffer-file-name)) "&"))
%  presentation (shell-command (concat "evince " (replace-regexp-in-string "tex$" "pdf"(buffer-file-name)) "&"))

%http://www.naney.org/diki/dk/hyperref.html
%日本語EUC系環境の時
\AtBeginDvi{\special{pdf:tounicode EUC-UCS2}}
%シフトJIS系環境の時
%\AtBeginDvi{\special{pdf:tounicode 90ms-RKSJ-UCS2}}

\title{東京エリアDebian勉強会}
\subtitle{第73回 2011年2月度}
\author{まえだこうへい mkouhei@debian.org\\IRC nick: mkouhei}
\date{2011年2月19日}
\logo{\includegraphics[width=8cm]{image200607/openlogo-light.eps}}

\begin{document}

\frame{\titlepage{}}

\section{}
\begin{frame}
 \frametitle{Agenda}
\begin{minipage}[t]{0.45\hsize}
  \begin{itemize}
  \item 注意事項
	\begin{itemize}
	 \item 飲酒禁止
	 \item 宗教禁止
	 \item 営利活動禁止
	\end{itemize}
  \item 最近あったDebian関連のイベント報告
	\begin{itemize}
	 \item 2010年1月勉強会
	\end{itemize}
 \end{itemize}
\end{minipage} 
\begin{minipage}[t]{0.45\hsize}
 \begin{itemize}
  \item Hadoopネタ
  \item Debian Games Team体験記
  \item　Squeeze リリースプレパーチー
  \item CACert認証会 \& GPGキーサイン
 \end{itemize}
\end{minipage}
\end{frame}

\begin{frame}
 \frametitle{前回}
\begin{minipage}[t]{0.45\hsize}
  \begin{itemize}
  \item 注意事項
	\begin{itemize}
	 \item 飲食禁止
	 \item 宗教禁止
	 \item 営利活動禁止
	\end{itemize}
  \end{itemize}
\end{minipage}
\begin{minipage}[t]{0.45\hsize}
\begin{itemize}
  \item 最近あったDebian関連のイベント報告
	\begin{itemize}
	 \item Debian勉強会アンケートシステム
	 \item kinect
	 \item CACert assurance(宴会で)
	\end{itemize}
 \end{itemize}
\end{minipage}
\end{frame}


\emtext{イベント報告}

\section{DWN quiz}
\begin{frame}{Debian 常識クイズ}

Debian の常識、もちろん知ってますよね?
知らないなんて恥ずかしくて、知らないとは言えないあんなことやこんなこと、
みんなで確認してみましょう。

今回の出題範囲は\url{debian-devel-announce@lists.debian.org} に投稿された
内容とDebian Project Newsからです。

\end{frame}

\subsection{問題}
%; whizzy-master ../debianmeetingresume201101.tex
% $B0J>e$N@_Dj$r$7$F$$$k$?$a!"$3$N%U%!%$%k$G(B M-x whizzytex $B$9$k$H!"(Bwhizzytex$B$,MxMQ$G$-$^$9!#(B
%
% $B$A$J$_$K!"%/%$%:$OJL%V%i%s%A$G:n@.$7!"$N$A$K%^!<%8$7$^$9!#5U$K%^!<%8$7(B
% $B$J$$$h$&$K$7$^$7$g$&!#(B
% (shell-command "git checkout quiz-prepare")

\santaku
{Debian 6.0 $B$,%j%j!<%9$5$l$?$N$O$$$D$+!)(B}
{2011/02/05}
{2011/02/06}
{2011/02/07}
{A}
{}

\santaku
{$B:#2s$N%j%j!<%9$G%5%]!<%H$5$l$J$/$J$C$?%"!<%-%F%/%A%c$O!)(B}
{m68k, ppc64}
{alpha, arm, hppa}
{blackfin, microblaze}
{B}
{}

\santaku
{www.debian.org $B$N%G%6%$%s$,JQ99$5$l$^$7$?!#2?G/F1$8%G%6%$%s$r;H$C$F$$$?$+!)(B}
{11$BG/(B}
{12$BG/(B}
{13$BG/(B}
{C}
{}

\santaku
{$B<!$N%j%j!<%9%3!<%I%M!<%`$O2?$+!)(B}
{rc}
{spell} % Mr.Spell
{wheezy}
{C}
{}

\santaku
{$B:#EY9T$o$l$k(BFTP-master$B2q5D$G5DO@$5$l$J$$M=Dj$O!)(B}
{$B%Q%C%1!<%8<+F0%5%$%s5!9=(B}
{$B%/%m%9%3%s%Q%$%k%5%]!<%H(B}
{md5sum$B$H(BshaXsum$B$r$H$C$Q$i$&(B}
{B}
{}

\santaku
{Lars Wirzenius$B$,Ds0F$7$?<!$N%j%j!<%9%4!<%kL\I8$O!)(B}
{UTF-8$B$GF0$+$J$$%Q%C%1!<%8$r$J$/$9(B}
{Debian$B$GF0$$$F$$$k?M9)1R@1$rBG$A>e$2$^$9(B}
{$BA4%Q%C%1!<%8$r(BLLVM$B$G%S%k%I$7$^$9(B}
{A}
{}


\emtext{Hadoopネタ}

\emtext{Debian Games Team体験記}

\emtext{prework}

{\footnotesize
 %; whizzy-master ../debianmeetingresume201101.tex
% $B0J>e$N@_Dj$r$7$F$$$k$?$a!"$3$N%U%!%$%k$G(B M-x whizzytex $B$9$k$H!"(Bwhizzytex$B$,MxMQ$G$-$^$9!#(B

\begin{prework}{ $B%-%?%O%i(B }

$B<B2H$K%j%j!<%9A0$N(BSqueeze$B$rCV$$$F$-$?$,!"FC$KLdBj$J$/;HMQ$5$l$F$$$kMM;R!#(B
$B2?$b@_Dj$7$J$/$F$b!"%5%9%Z%s%I!&%O%$%P%M!<%7%g%s$G$-$?$h$&$J!"%W%j%s%?$N@_Dj$K>/!9G:$s$@5-21$,$"$k$,!">\:Y$OK:$l$?!#(B
$B@N$KHf$Y$k$H3JCJ$K3Z$K$J$j$^$7$?!#(B
\end{prework}

\begin{prework}{ $B5HLn(B(yy\_y\_ja\_jp) }

poppler-data,cmap-adobe-*,gs-cjk-resource $B$,(B main $B$KF~$C$?$3$H$G$9$M!$(BThanks to Adobe$B!%(B
\end{prework}

\begin{prework}{ d123.itouch.com(dai) }

 \begin{itemize}
  \item $B1'Ch$,%F!<%^$G$+$C$3$$$$!*(B
  \item $B?7$7$/(BDebian$B$r3P$($k$3$H$K$J$C$?$N$G!"CzEY$$$$5!2q$G$9!#(BDebian$B=i(B
	$B?4<T!*0lG/@8$G$9!*(B
 \end{itemize}

 2$BG/$[$IA0$KEl5~$NK?@lLg3X9;$G3+:E$5$l$?5.2q$N%V!<%9$G$4CzG+$K(B
 $B$4>R2p$$$?$@$$$?$3$H$,$"$j$^$9!#$9$P$i$7$$2q$r3+:E$$$?$@$-(B
 $B$"$j$,$H$&$4$6$$$^$9!#:#2s=i$N;22C$G$9!#$I$&$>$h$m$7$/$*4j$$$7$^$9!#(B
\end{prework}

\begin{prework}{ hakase }

$B$9$$$^$;$s!"(BSqueeze$B$NL>A0;O$a$FJ9$-$^$7$?!#(B
$BJY6/$7$^$9!#(B
\end{prework}

\begin{prework}{ henrich }

$B$9$$$^$;$s!"FCCJ8D?ME*$K$&$l$7$$$H$3$m$,8+Ev$?$j$^$;$s$G$7$?!D(B
\end{prework}

\begin{prework}{ $B;3ED(B }

\preworksection{$B$&$l$7$$$3$H(B}
\begin{itemize}
 \item /bin/sh$B$,(Bdash$B$K$J$C$?!*%9%/%j%W%H5/F0$,L@$i$+$K7Z$$(B
 \item $B0MB84X78%Y!<%9$G5/F0=hM}$,$5$l$k$h$&$K$J$C$?!JAa$/$J$C$?!K(B
 \item $B%Q%C%1!<%8$,$5$i$K=<<B$7$?!J5$$E$$$F$J$+$C$?$@$1$+$b$7$l$J$$$b$N$N!"!V$3$s$J$N$^$GF~$C$F$k!WEY$,>e>:!K(B
 \item $B%j%j!<%9%9%1%8%e!<%k$,$7$C$+$j$7$F$-$?(B
\end{itemize}

\preworksection{$BJQ$o$C$?$3$H!J5$$K$J$k$3$H!K(B}
\begin{itemize}
 \item kFreeBSD$B$,F~$C$?!*$9$0;H$&$o$1$G$O$J$$$1$I!"5$$K$J$k(B
 \item arm(eb)$B$,Mn$C$3$A$?$1$I!"@$$NCfBgBN(Barmel$B$@$+$i(BOK$B$J$N$+$J!)(B
 \item Debian Live$B!#(BLiveCD$B$+$i(Bnetboot$B$^$G%+%P!<$H$$$&$3$H$GMW%A%'%C%/(B
\end{itemize}
$B%j%j!<%9$NA0$+$i<B:]$K$O;H$($F$$$?$N$G!V@hF|$+$i$$$-$J$j$&$l$7$/$J$C$?!W;v$O<B$O$J$$$b$N$N!"Ce<B$J%j%j!<%9$H$$$&46A[$G$9!#(B
\end{prework}

\begin{prework}{ dictoss($B?yK\!!E5=<(B) }

$B5/F0$,Aa$/$J$C$?!#(B
Debian GNU/kFreeBSD$B$NL>$,9-$^$C$F$-$?!#(B
\end{prework}

\begin{prework}{ yamamoto }

\preworksection{$B$&$l$7$$$3$H(B}
Squeeze $B$@$+$i!"$H$$$&$o$1$G$O$J$$$G$9$,!"(Bdebian-multimedia $B$N(B mythtv $B$,(B
 0.24 $B7O$K$J$C$F!"(Bmythtv-backend $B$,FMA3;`$9$k$3$H$,L5$/$J$j$^$7$?!#(B

\preworksection{$BJQ$o$C$?$3$H(B}
KDE $B$r>oMQ$7$F$^$9$,!"(BKlipper $B$N5sF0$,JQ$o$C$?!#$J$s$+;~!9!"$o$6$o$6(B Klipper $B$N%A%'%C%/$r$D$1$J$$$H!"Cf%\%?%s$GA*BrJ8;zNs$r%Z!<%9%H$G$-$J$$!#(B
$B860xD4::$O$^$@$@$1$I!D!#(B
$B!V%/%j%C%W%\!<%I$N%"%/%7%g%s$rM-8z$K$9$k!W$"$?$j$,$"$d$7$$(B!?
\end{prework}

\begin{prework}{ $BLnEg!!5.1Q(B }

$B$&$O$C!"IaCJ$+$i(Bunstable(sid)$B$7$+;H$C$F$J$$(B...

$B<h$j5^$.(Blenny$B$H$N0c$$$G4r$7$$;v!#(B
\begin{enumerate}
 \item kernel$B$,(B2.6.32$B7ONs$K$J$C$?!#(B(kvm$B$N%P!<%8%g%s$,>e$,$C$?!*(B)
 \item libvirt$B$,(B0.8.3$B$K$J$C$?!#!J$$$m$s$J%Q%i%a%?$,$$$m$$$m$D$1J|Bj$K(B...)
 \item gnome$B$,(B2.30$B7ONs$K$J$C$?!#(Bfreedesktop-sound-theme$B$,I8=`AuHw$5$l$?$N$G2;$,$K$.$d$+$K$J$C$?!#(B
 \item upstart$B$,Ek:\$5$l$?!#(B(sid$B;H$C$F$$$k$N$G$9$,!"(Bupstart$B$K$J$C$F$J$/(B
       $B$F$,$C$/$j!#:#$+$iJQ$($F$_$^$9(B...)
\end{enumerate}
\end{prework}

\begin{prework}{ taitioooo }

$B=i$a$F(BDebian$B$K?($C$?$N$,(Bsqueeze$B$G$7$?!#(B
$B:#$^$G(BCentos$B;H$C$F$$$^$7$?!#(B
$B=i$a$F$N0u>]$O!"!V$"!"%m%1%C%H$,$"$k!#!W$G$7$?!#(B

$B$"$H!"(Bapt-get$B$N@_Dj$,4JC1$G$h$+$C$?$G$9!#!J(Byum$B$HHf$Y$F$G$9$,!K(B
\end{prework}

\begin{prework}{ Hasegawa }

$B%$%s%9%H!<%k$G(Broot$B$G%m%0%$%s$9$k$+J9$$$F$-$?$j(B
$B@_Dj9`L\$,A}$($?$h$&$@$7!"$+$J$jJQ$o$C$?0u>]$r;}$C$?(B
64bit/32bit$B$,A*Br$G$-$k%$%s%9%H!<%i!<$K$J$C$?(B

\end{prework}

\begin{prework}{ ayakomuro }

$B$^$@%"%C%W%0%l!<%I$7$F$J$$$N$G!"=5Kv$K%"%C%W%0%l!<%I$7$?$i>\:Y=q$-$^$9(Bm(..)m
\end{prework}

\begin{prework}{$B$^$($@$3$&$X$$(B}

\preworksection{$B$&$l$7$$$3$H(B}
\begin{itemize}
 \item Squeeze$B$,%j%j!<%9$5$l$?$3$H$G(BSid$B$N3+H/$,:F3+$5$l$?$N$G!"(BCouchApp$B$d(B
       xserver-xorg-input-multitouch$B$,(BDebian$B%Q%C%1!<%8$K$J$C$?!#(B
 \item Sarge$B$G2TF/Cf$NAH9g%5!<%P$r%O!<%I%&%'%"!u%7%9%F%`O75`2=$K$h$j!"(B
       KVM/Squeeze/HP MicroServer$B$G9=C[Cf$J$N$G!"%j%j!<%9$5$l$F$A$g$&$I(B
       $BNI$+$C$?!#(B
 \item $B:#$^$G(BSid$B$G$7$+;H$($J$+$C$?%Q%C%1!<%8$,(Bstable$B$G;H$($k$h$&$K$J$C$?(B
       $B$b$N$,A}$($?!#(B(lxc, scala, xz-utils, libapache-mod-security($B$3$l$OI|3h$7(B
       $B$?!"$H$$$&$N$,@5$7$$$+(B))
 \item $B5/F0%9%/%j%W%H$,(Bbash$B$+$i(Bdash$B$KJQ$o$C$?!#(BVM$B$G$b5/F0$,B.$/$J$C$?5$(B
       $B$,$9$k!#(B($BBN46B.EY!#(B)
\end{itemize}
\preworksection{$BJQ$o$C$?$3$H(B}
\begin{itemize}
 \item libvirt$B$,(B0.8.3$B$K$J$C$?!#(Bnwfilter$B$N5!G=$,F~$C$?(B
       (/etc/libvirt/nwfilter/*.xml)$B!#(BNetfilter$B$N%k!<%k$O$9$Y$F<+J,$G=q(B
       $B$$$F$$$k$N$G!"<+J,$K$OI,MW$J$$$1$l$I!"IaCJ=q$+$J$$(BNetfilter$B$N%k!<(B
       $B%k$r=q$+$J$$?M$K<h$C$F$ONI$$$N$G$O!#(B
 \item $B%/%i%&%I4XO"$N%D!<%k$,%Q%C%1!<%8$GA}$($?!#;~Be$G$9$M!<!#(B
\end{itemize}
\end{prework}



}

\begin{frame}{インストールした機器}
\begin{table}[ht]
 \begin{center}
  \begin{tabular}{|l|l|}
   \hline
   メーカー & 機種名 \\
   \hline
   HP & MicroServer  \\
   \hline
  \end{tabular}
 \end{center}
\end{table}
\end{frame}

\begin{frame}{インストールして困ったこと}

\begin{table}[ht]
 \begin{center}
  \begin{tabular}{|p{7em}|p{4em}|p{7em}|p{4em}|}
   \hline
   困ったこと & 解決方法 & 宿題 & 担当 \\
   \hline
   プリンタの設定に少々悩んだ & 忘れた & 思いだしてフィードバック & キタハラ\\
   \hline
  \end{tabular}
 \end{center}
\end{table}

\end{frame}

\begin{frame}{改善した方が良いと思うこと。}

\begin{table}[ht]
 \begin{center}
  \begin{tabular}{|p{8em}|p{11em}|p{3em}|}
   \hline
   改善した方良いと思ったこと & 宿題 & 担当 \\
   \hline
   libvirt のnwfilterの設定によって、自分で設定しているiptablesのルール
   が上書きされる & nwfilterにルールを追加、又は全て削除してス
	   クリプトに集約。前者はルールの順番が不明なので仕組みを理解 & まえだ\\
   \hline
   Klipperで中ボタンでのペーストが不便になった& 「クリップ
	   ボードのアクションを有効にする」で解決するか確認。NGならば改善案をBTS & 山本\\
   \hline
   心残りがあること & すっきり解決する & 岩松 \\
   \hline
   d-iでLVMを一度設定すると削除不可　& シェルで\{lv,vg,pv\}remove
       すれば削除可。メニューから削除できた方が便利 & まえだ \\
   \hline
  \end{tabular}
 \end{center}
\end{table}

\end{frame}

\begin{frame}{今後のイベント}
 
\begin{itemize}
 \item 3月 OSC
 \item CACert keysigning
 \item OSC Kobe
 \item 3月第三土曜 Debian Coffee?
 \item 4月　PS Moveネタ? デジタル放送取り込み? Debian Pod cast? 100台
       Squeezeアップグレード(吐血)体験記? 
\end{itemize}
\end{frame}

\begin{frame}{今日の宴会場所}

大久保駅周辺。

\end{frame}

\end{document}

;;; Local Variables: ***
;;; outline-regexp: "\\([ 	]*\\\\\\(documentstyle\\|documentclass\\|emtext\\|section\\|begin{frame}\\)\\*?[ 	]*[[{]\\|[]+\\)" ***
;;; End: ***
